\section{Related Work}
\label{sec:related_work}

\subsection{Computer Control and Game Environments}
\label{sec:related_work_computer_control}

Benchmarks for agents built on foundation models mainly come from two lines.
The first studies software and general computer control, including web,
mobile, desktop, and real computer tasks~\citep{liu2024agentbench,zhou2023webarena,tan2024cradle}.
Recent benchmarks such as ScenDroid and GameWorld/V-MAGE further sharpen this line~\citep{scendroid2026,ouyang2026gameworld,zheng2025vmage}.
These settings are useful for tool use and interface grounding, but they only
partly capture the navigation, scene changes, and recovery pressure of dynamic
open-world games.
ALFWorld aligns text interaction with embodied household environments, while
MineDojo provides an open-ended Minecraft benchmark with internet-scale
knowledge~\citep{shridhar2021alfworld,fan2022minedojo}. Voyager shows that
large models can support long-running exploration in Minecraft-like
worlds~\citep{wang2024voyager}. TextAtari extends this direction to 100K-frame
decision-making and makes the trade-off between expensive reasoning and
reactive execution explicit~\citep{textatari2025}.
\stardojo{}~\citep{tan2025stardojo} lies between these traditions: it keeps
open-world dynamics but requires control without a privileged game-state API,
using screen-derived observations and a shared executable action interface,
making low-level reliability and control scheduling central.

\subsection{LLM/VLM Agents for Computer Tasks and Games}
\label{sec:related_work_llm_vlm_agents}

LLM and VLM agents have made strong progress on multi-step computer
interaction. ReAct established the observe-reason-act loop~\citep{yao2023react}.
Later agents extended this idea to web navigation, app use, and general
computer control~\citep{zhou2023webarena,deng2023mind2web,tan2024cradle}.
In games, Voyager demonstrates strong large-model planning, while TextAtari
shows how long-horizon tasks expose the tension between costly reasoning and
reactive execution~\citep{wang2024voyager,textatari2025}.
DEPS, Lumine, and SIMA extend this line with open-world control~\citep{wang2023deps,tan2025lumine,abiraad2024scaling},
while JARVIS-1 adds memory-augmented multimodal control~\citep{wang2023jarvis1}.

\paragraph{Scope of comparison.}
Recent open-world game-agent systems and benchmarks are important context but
are not always direct same-table baselines for \method{}. NitroGen trains a
vision-action foundation model via large-scale behavior cloning on 40,000 hours
of gameplay across more than 1,000 games~\citep{magne2026nitrogen}, whereas
\method{} studies test-time controller allocation without additional agent
training. Lumine targets real-time 3D open-world control with an end-to-end VLM
agent~\citep{tan2025lumine}, while \stardojo{} Lite-100 uses benchmark tasks,
shared runner feedback, and a fixed executable action interface. GameWorld
standardizes verifiable evaluation across 34 browser games and 170 tasks, and
V-MAGE provides a vision-centric protocol~\citep{ouyang2026gameworld,zheng2025vmage}.
Because these works differ in training, environments, action interfaces, and
verification signals, we discuss them as related positioning rather than direct
baselines.

However, many game agents still rely on privileged observations, task
scaffolding, or in-environment training. \method{} instead focuses on a setting
with no training and no privileged game-state API: the Reactive Controller must
remain responsive under screen-grounded interaction, while strategic reasoning,
used only when needed, handles drift, recovery, and task-level redirection.

\subsection{Dual Controller Frameworks and Memory for Long-Horizon Agents}
\label{sec:related_work_dual_controller_memory}

Fast-slow and hierarchical control manage the trade-off between costly reasoning
and real-time execution. SayCan combines language planning with affordances, Inner Monologue
uses closed-loop language feedback, and SwiftSage explores fast-slow
cooperation~\citep{ahn2022saycan,huang2022innermonologue,lin2023swiftsage}.
ADaPT and AdaPlanner further study on-demand decomposition and
feedback-adaptive planning~\citep{prasad2024adapt,sun2023adaplanner}.
Search methods, including Tree of Thoughts, RAP, and LATS, show the value of
deliberate exploration~\citep{yao2023treeofthoughts,hao2023rap,zhou2024lats}.
Reflexion further highlights feedback and retry for recovery~\citep{shinn2023reflexion}.
Many prior agents remain planning-intensive or rely on fixed triggering. In
contrast, \method{} uses multimodal event triggering, local autonomous override,
and graded recovery so that strategic reasoning is called only when needed.

Memory is equally important for long-horizon behavior. Generative Agents,
Reflexion, and MemGPT show that retrieval and reflection can help agents go
beyond a single prompt window~\citep{park2023generativeagents,shinn2023reflexion,packer2023memgpt}.
HiAgent and Mem0 further study hierarchical working memory and scalable
long-term memory~\citep{hu2025hiagent,chhikara2025mem0}. Retrieval-Augmented
Generation (RAG), GraphRAG,
and HippoRAG show how external stores and graph retrieval can extend model
context~\citep{lewis2020rag,edge2024graphrag,gutierrez2024hipporag}.
Recent agents without privileged environment APIs also use state-action graphs
for exploration and delayed reward planning~\citep{tang2025experiencedriven}.
AgentOdyssey studies open-ended text games with test-time continual learning,
highlighting exploration, episodic memory, and planning cost~\citep{agentodyssey2026}.
Our memory design follows this direction but separates roles: \samb{} supports
fast local reuse, whereas SA-KG provides evidence for strategic checking and
replanning.
